% !TeX root = ../../Book1.tex
% 纲要标题：### 15.2 代数闭包
% 纲要位置：计划纲要.md，第 536 行。

\section{代数闭包}
\label{b1:sec:1502}


% 纲要内容（仅作写作计划，尚非教材正文）：
%
% * 代数闭包的存在；有限多条根关系的相容性与递增并域
% * 一般嵌入延拓与代数闭包在基域同构意义下的唯一性
% * 选择原理在证明中的位置
%

一个分裂域只处理指定多项式。若要在同一个域中讨论基域上的所有代数方程，可先构造代数扩张，使 \(K[x]\) 中的每个非常数多项式至少有一个根。但得到新域 \(F_1\) 后，\(F_1[x]\) 中还可能出现系数不全属于 \(K\) 的多项式，刚才的构造并未为这些方程安排根。因此对新域继续添根，得到递增的域列 \(F_0=K\subseteq F_1\subseteq F_2\subseteq\cdots\)。利用每个多项式只有有限多个系数，可以证明这些域的并处理了其中全部的多项式方程。

\subsection{代数闭包的存在性}
\label{b1:subsec:1502:01}

\begin{definition}\label{b1:def:algebraic-closure}
设 \(\overline K/K\) 为域扩张。如果 \(\overline K/K\) 是代数扩张，并且 \(\overline K\) 是代数闭域，即 \(\overline K[x]\) 中每个非常数多项式都在 \(\overline K\) 中有根，则称 \(\overline K\) 为 \(K\) 的\term[daishu-bibao]{代数闭包}[algebraic closure]。
\end{definition}
\symidx{\(\overline K\)：\(K\) 的一个代数闭包}

要使域 \(F\) 上的每个非常数多项式都有根，可以为每个首一非常数多项式 \(f\) 设置一个形式符号 \(X_f\)，再在多项式环中施加关系 \(f(X_f)=0\)。若这些关系在一个非零商环中同时成立，\(X_f\) 的像便成为 \(f\) 的根。下面还要使这个商环成为域，并保证所得扩张代数。

\begin{lemma}\label{b1:lem:adjoin-all-roots}
对任意域 \(F\)，存在代数扩张 \(F'/F\)，使 \(F[x]\) 中每个非常数多项式在 \(F'\) 中至少有一个根。
\end{lemma}
\begin{proof}
为每个首一非常数多项式 \(f\in F[x]\) 设置一个不定元 \(X_f\)，令 \(R=F[X_f\mid f]\)。每个环元素只涉及有限多个不定元。令 \(I\) 为全部 \(f(X_f)\) 生成的理想；在 \(R/I\) 中，这些关系便同时成立。

我们先证 \(I\neq R\)。关键是任何有限多个这样的关系都能在某个扩域中同时实现。如果 \(1\in I\)，按生成理想的定义，就有有限表达式
\[
1=\sum_{i=1}^r h_i f_i(X_{f_i}),\qquad h_i\in R,
\]
其中可将重复的 \(f_i\) 合并。由\cref{b1:prop:splitting-field-exists}逐次为 \(f_1,\ldots,f_r\) 添入根，得到某个扩张域 \(E/F\)，在其中同时选定 \(f_i\) 的根 \(a_i\)。把 \(X_{f_i}\) 送到 \(a_i\)，其余不定元送到零，就得到保持 \(F\) 的环同态 \(R\rightarrow E\)。余变量取零只是为了给整个多项式环定义同态，并不要求这次赋值满足它们各自对应的添根关系；上式只需所列的有限条关系成立。上式的右边映为零，左边仍为一，矛盾。因此有限子系统的相容性排除了 \(1\in I\)。

由\cref{b1:prop:maximal-ideal-exists}，真理想 \(I\) 包含于一个极大理想 \(\mathfrak m\)。令 \(F'=R/\mathfrak m\)，则它是域；\(F\cap\mathfrak m=0\)，因为非零常数可逆，所以 \(F\) 嵌入 \(F'\)。记 \(X_f\) 的像为 \(\overline X_f\)，则 \(f(\overline X_f)=0\)，所以每个 \(\overline X_f\) 都在 \(F\) 上代数。任意 \(a\in F'\) 都是有限多个 \(\overline X_{f_1},\ldots,\overline X_{f_s}\) 的多项式表达，因而属于
\[
F(\overline X_{f_1},\ldots,\overline X_{f_s}).
\]
由\cref{b1:lem:finite-algebraic-generators}，这个中间域在 \(F\) 上有限，再由\cref{b1:lem:finite-dimensional-algebraic}知 \(a\) 在 \(F\) 上代数。因此 \(F'/F\) 是代数扩张。任意非常数多项式除以首项系数便成为首一多项式，根不改变，故所有非常数多项式都有根。
\end{proof}

\begin{theorem}\label{b1:thm:algebraic-closure}
每个域 \(K\) 都存在代数闭包。
\end{theorem}
\begin{proof}
从 \(F_0=K\) 出发，反复应用\cref{b1:lem:adjoin-all-roots}构造代数扩张 \(F_{n+1}/F_n\)，使 \(F_n\) 上每个非常数多项式在 \(F_{n+1}\) 中有根。将各步嵌入认同为包含，令 \(\Omega=\bigcup_{n\geq0}F_n\)。任意两个元素都落在某个共同的 \(F_n\) 中，其差、积以及非零元素的逆仍在该域内，因此 \(\Omega\) 是域。反复应用\cref{b1:thm:algebraic-tower}可知每个 \(F_n/K\) 代数。任意 \(a\in\Omega\) 属于某个 \(F_n\)，所以在 \(K\) 上代数；这证明了 \(\Omega/K\) 代数。

若 \(g=\sum_{i=0}^d c_ix^i\in\Omega[x]\) 非常数，每个系数 \(c_i\) 属于某个 \(F_{n_i}\)。系数只有有限多个，取 \(n=\max\{n_0,\ldots,n_d\}\)，并利用域列递增，就有所有 \(c_i\in F_n\)。所以 \(g\in F_n[x]\)，在 \(F_{n+1}\subseteq\Omega\) 中有根。故 \(\Omega\) 代数闭，满足代数闭包的两项要求。
\end{proof}
这项构造不要求基域可数。每一步同时为该域的全部多项式安排不定元；迭代后的有限系数论证直接证明并域代数闭，不需要先判定第一步所得域是否已代数闭。“可数次迭代”也不意味着所得代数闭包是可数集。

\subsection{一般嵌入延拓与代数闭包的唯一性}
\label{b1:subsec:1502:02}

\begin{theorem}[嵌入延拓]\label{b1:thm:embedding-extension}
设 \(L/K\) 为任意代数扩张，\(\Omega\) 为代数闭域，\(\sigma:K\rightarrow\Omega\) 为域嵌入。则存在嵌入 \(\widetilde\sigma:L\rightarrow\Omega\) 延拓 \(\sigma\)。
\end{theorem}
\begin{proof}
只要一个部分嵌入的定义域还没有覆盖 \(L\)，就尝试再添入一个元素并延拓。Zorn 引理给出极大部分嵌入，而代数性将保证其定义域不能遗漏任何元素。

令 \(\mathcal P\) 为所有 \((E,\tau)\) 组成的集合，其中 \(E\) 是 \(L/K\) 的中间域，\(\tau:E\hookrightarrow\Omega\) 是满足 \(\tau|_K=\sigma\) 的域嵌入。规定 \((E,\tau)\le(E',\tau')\) 当且仅当 \(E\subseteq E'\) 且 \(\tau'|_E=\tau\)。初始对 \((K,\sigma)\) 使 \(\mathcal P\) 非空。

对一条非空链 \(\mathcal C\)，将其定义域取并，得到中间域 \(E_{\mathcal C}\)。链中任意两个定义域按包含关系可比，较大域上的映射限制到较小域后就是原映射，所以两个映射在重叠处绝不冲突。于是可定义 \(\tau_{\mathcal C}(a)=\tau(a)\)，其中任选一个包含 \(a\) 的链成员 \((E,\tau)\)，所得值与选择无关。任意有限多个元素都可放进同一个链成员的定义域，在那里验证 \(\tau_{\mathcal C}\) 保持域运算并且单射，就知它是延拓 \(\sigma\) 的嵌入。故 \((E_{\mathcal C},\tau_{\mathcal C})\) 给出上界；空链以 \((K,\sigma)\) 为上界。Zorn 引理（\cref{b1:thm:A-zorn}）给出极大对 \((E,\tau)\)。

若 \(E\neq L\)，取 \(a\in L\setminus E\)。因为 \(L/K\) 代数，\(a\) 满足某个非零的 \(K\) 系数多项式方程。\(K\subseteq E\)，所以同一个方程也是非零的 \(E\) 系数方程，\(a\) 因而在 \(E\) 上代数。其最小多项式经系数同构 \(\tau:E\rightarrow\tau(E)\) 变为 \(\tau(E)\) 上的不可约多项式。\(\Omega\) 代数闭，故该多项式在 \(\Omega\) 中有根。\cref{b1:lem:simple-embedding-extension}便把 \(\tau\) 延拓到 \(E(a)\)，与极大性矛盾。所以 \(E=L\)，得到所需延拓。
\end{proof}

\begin{proposition}\label{b1:prop:algebraic-closure-unique}
若 \(\Omega/K\)、\(\Omega'/K\) 都是代数闭包，则存在保持 \(K\) 的同构 \(\Omega\rightarrow\Omega'\)。
\end{proposition}
\begin{proof}
对代数扩张 \(\Omega/K\) 和嵌入 \(K\hookrightarrow\Omega'\) 应用\cref{b1:thm:embedding-extension}，得到保持 \(K\) 的嵌入 \(\tau:\Omega\rightarrow\Omega'\)。为了证明像域 \(\tau(\Omega)\) 代数闭，任取其中的非常数多项式
\[
g(x)=\sum_{i=0}^d\tau(a_i)x^i,
\qquad a_i\in\Omega,
\]
将系数拉回，得到 \(g^\tau(x)=\sum_{i=0}^d a_ix^i\in\Omega[x]\)。因 \(\tau\) 单射，这个多项式仍非常数。\(\Omega\) 代数闭，故有 \(c\in\Omega\) 使 \(g^\tau(c)=0\)。映回去便有
\[
g(\tau(c))=\tau\bigl(g^\tau(c)\bigr)=0,
\]
所以 \(g\) 在像域中有根。

任意 \(b\in\Omega'\) 在 \(K\) 上代数。由于 \(\tau\) 保持 \(K\)，有 \(K\subseteq\tau(\Omega)\)；原来零化 \(b\) 的非零 \(K\) 系数多项式仍是非零的 \(\tau(\Omega)\) 系数多项式，所以扩大基域以后 \(b\) 仍代数。其在像域上的不可约最小多项式只能为一次，因为像域代数闭。于是 \(b\in\tau(\Omega)\)，所以 \(\tau\) 满射。这里利用了“代数闭域没有非平凡代数扩张”，而非错误地认为无限集上的单射自动为满射。
\end{proof}
同构通常并不唯一；选择不同共轭根会给出不同延拓。后文写“固定一个代数闭包”时，是选定共同环境，后续的嵌入与根都在这个选定的域中讨论。

\subsection{证明中的选择原理}
\label{b1:subsec:1502:03}

\cref{b1:lem:adjoin-all-roots}由真理想取得极大理想，调用了先前用 Zorn 引理证明的极大理想存在性；\cref{b1:thm:embedding-extension}则直接对部分嵌入使用 Zorn 引理。

\ProseParagraphBreak
对链取并时必须检查两件不同的事：定义域之并对域运算封闭；映射在重叠处给出同一个值。只有在偏序要求“延拓”而不只是“定义域更大”时，第二件事才得到保证。极大部分嵌入的论证还依赖 \(L/K\) 代数：若剩余元素超越，不能再以最小多项式及其根作为延拓依据。

\ProseParagraphBreak
本书在这些一般存在性论证中采用通常的选择公理。算法性问题另作处理：抽象地得到一个极大理想或代数闭包，并未给出可计算的元素编码、多项式分解程序或根的选择算法。具体基域上的有效计算，须另加适当表示方法。
